\subsection{Does Language Shape Thought?}
\label{sec:language-and-thought}
\slidesrc{Pragmatics: 11--49}

``Thoughts shape language\ldots\ but\ldots\ does language shape thought?'' If pragmatics
studies how context and intentions shape what we mean, the reverse question is natural:
does the language we speak shape how we think?

\subsubsection{Newspeak and the Sapir--Whorf hypothesis}
\slidesrc{Pragmatics: 13--16}

In George Orwell's novel \emph{Nineteen Eighty-Four} (1949), the regime designs a new
language, called \textbf{Newspeak}, to control thought:

\begin{quote}\small
``Don't you see that the whole aim of Newspeak is to narrow the range of thought? In the
end we shall make thought\-crime literally impossible, because there will be no words in
which to express it. Every concept that can ever be needed will be expressed by exactly
one word, with its meaning rigidly defined and all its subsidiary meanings rubbed out and
forgotten. [\ldots] Every year fewer and fewer words, and the range of consciousness
always a little smaller. [\ldots] The Revolution will be complete when the language is
perfect.''
\end{quote}

\begin{nbdef}{Linguistic relativity (Sapir--Whorf hypothesis)}
The hypothesis of \textbf{linguistic relativity} (also \textbf{Sapir--Whorf
hypothesis}, or Whorfianism, after Edward Sapir, 1884--1939, and Benjamin Lee Whorf,
1897--1941) holds that the structure of a language affects its speakers' world view or
cognition. It has two versions:
\begin{itemize}[nosep]
  \item the \textbf{strong} version, \textbf{linguistic determinism}: language
        \emph{determines} thought, linguistic categories limit and determine cognitive
        categories; language entirely determines the range of cognitive processes;
  \item the \textbf{weak} version, \textbf{linguistic influence}: linguistic categories
        and usage only \emph{influence} thought and decisions; language provides
        constraints in some areas of cognition, but is by no means determinative.
\end{itemize}
\end{nbdef}

\begin{nbfact}{what we know today}
Linguistic \textbf{determinism} (the strong version) is now generally agreed to be
\textbf{false}: we can think things for which our language has no word (otherwise new
words could never be invented). Research on the \textbf{weaker} forms has produced
\textbf{positive empirical evidence}: language does influence some areas of cognition.
Newspeak is the strong version, and that is why it is fiction.
\end{nbfact}

\subsubsection{How languages and cultures shape the ways we think}
\slidesrc{Pragmatics: 17--30}

The rest of this part follows a talk by \textbf{Lera Boroditsky} (Associate Professor of
Cognitive Science at UC San Diego, previously at MIT and Stanford, who studies the
relationships between mind, world and language), \emph{How languages and cultures shape
the ways we think}.

\paragraph{Languages force us to say different things.}
Take the English sentence \emph{``Trump read Ta-Nehisi's latest book''}. Depending on the
language, to say it you \emph{must} decide:
\begin{itemize}[nosep]
  \item did this happen in the past? (tense: Mandarin does not require it, English does);
  \item Donald or Melania? (in Russian the past-tense verb changes with the gender of the
        subject);
  \item was the action completed? (aspect: Russian forces a choice);
  \item did you witness this with your own eyes? (evidentiality: in Turkish, for
        example, the verb marks whether you saw it or heard about it).
\end{itemize}
Similarly, \emph{``the 7th blue house to the left of the big tree''}: seventh? blue, light
or dark (Russian has no single word for ``blue'')? left of, or north-northwest of, or
``volcano-ward'' of (some languages use absolute directions)? and which tree: banyan, red
cypress, eucalyptus? Each language makes its speakers pay attention to different
aspects of the world, all the time.

\paragraph{How many languages in the world?}
According to Ethnologue: 7\,168 languages in use (October 2023), 7\,164 (October 2024),
7\,159 (November 2025). The number is constantly in flux, because we keep learning about
languages and because languages themselves change. Roughly \textbf{44\%} of all languages
are \textbf{endangered}, often with fewer than 1\,000 users left, while the 20 largest
languages are the native tongue of more than 3.6--3.7 billion people: about 0.3\% of the
languages account for nearly half of the world's population. Another source is the
\emph{World Atlas of Language Structures} (WALS), which describes the structural
features of a sample of languages.

\paragraph{Navajo verbs of eating.}
Navajo has no single verb ``to eat'': the verb stem changes with the \emph{kind} of food
and the \emph{manner} of eating.

\begin{center}
\small
\begin{tabular}{@{}r l l l@{}}
\toprule
 & \textbf{stem} & \textbf{manner / food of consumption} & \textbf{example edibles} \\
\midrule
i.    & -y\'{\i}\'{\i}h & consume/chew \textsc{unspecified} & anything \\
ii.   & -'aa\l   & consume/chew \textsc{hard, compact} & corn, peyote, candy \\
iii.  & -chozh   & consume/chew \textsc{leafy} & lettuce, hay \\
iv.   & -gha\l   & consume/chew \textsc{meat} & meat \\
v.    & -k\'e\'ed & consume/chew \textsc{round} & bun, melon \\
vi.   & -ts'\'e\'eh & consume/chew \textsc{mushy matter} & mush, jello \\
vii.  & -dee\l   & consume/chew \textsc{plural object} & berries, eggs \\
viii. & -wol     & consume/chew \textsc{marrow} & marrow \\
ix.   & -joo\l   & consume \textsc{non-compact matter} & cotton candy \\
x.    & -t\l eeh & consume \textsc{mushy matter} & ice cream \\
xi.   & -m\'a\'a' & consume/devour by \textsc{gulping/bolting} & single piece of food \\
xii.  & -'\'a\'ah & consume \textsc{solid, round object} dunked in liquid & moistened
  bread, cake \\
xiii. & -dl\k{a} & consume/drink \textsc{unspecified} & soup, thin gruel \\
xiv.  & -kaah    & consume/drink from \textsc{open container} & milk from glass \\
xv.   & -t'aah   & consume/drink from \textsc{closed container} & milk from bottle \\
\bottomrule
\end{tabular}
\captionof{table}{Navajo verb stems of consumption (some diacritics simplified).}
\end{center}

\paragraph{Learning languages.} Quotes on the slides: ``To have a second language is to
have a second soul'' (Charlemagne); ``A man who knows four languages is worth four men''
(Charles V); ``I speak English to my accountants, French to my ambassadors, Italian to my
mistress, Latin to my God, and German to my horse'' (Frederick the Great). And the
opposite camp: ``I hate relativism more than I hate anything else, excepting, maybe,
fiberglass powerboats'' (\textbf{Jerry Fodor}, 1935--2017, one of the principal
philosophers of mind of the late twentieth century, very influential on cognitive
science, and a strong critic of linguistic relativity).

\paragraph{Why start from time?}
The most common nouns in English (lists exist also for Italian and French) are, in order:
\textbf{time}, year, people, way, day. Time is the most talked-about thing, so it is a
good place to look for differences between languages.

\subsubsection{Space and time}
\slidesrc{Pragmatics: 31--38}

We cannot see or touch time, so all languages talk about it using \textbf{space}; but
they use space differently.
\begin{itemize}
  \item \textbf{Writing direction}. English speakers arrange a temporal sequence (e.g.\
        pictures of a growing bear, as in the old Nestlé logo) from \emph{left to right};
        speakers of Hebrew or Arabic, who write right to left, arrange it \emph{right to
        left} (Tversky et al., 1991; Fuhrman \& Boroditsky, 2010).
  \item \textbf{Front and back}. For English speakers the future is \emph{ahead} and the
        past \emph{behind}. In \textbf{Aymara} (Andes) the past is \emph{in front} (it
        is known, you can ``see'' it) and the future \emph{behind}: speakers gesture
        forward when talking about the past (Núñez \& Sweetser, 2006).
  \item \textbf{Up and down}. In Mandarin, last week is the ``up'' week and next week the
        ``down'' week: time is also \emph{vertical}. So time can be laid on three axes:
        \textbf{transverse} (left--right), \textbf{sagittal} (front--back),
        \textbf{vertical} (up--down) (Boroditsky et al., 2010; Fuhrman et al., 2011).
  \item \textbf{Absolute directions}. In some Aboriginal Australian communities
        (Kuuk Thaayorre, Pormpuraaw) people do not use ``left'' and ``right'' but
        cardinal directions (``there is an ant on your south-west leg''). Given cards
        showing a man at different ages, they arrange them \textbf{from east to west},
        whichever way they are facing: left to right when facing south, towards their
        body when facing east (Boroditsky \& Gaby, 2010).
  \item \textbf{Uphill and downhill}. The Yupno of Papua New Guinea place the past
        downhill and the future uphill, following the slope of their valley (Núñez et
        al., 2012).
\end{itemize}

\begin{nbfact}{is language the cause?}
How do we know that language \emph{causes} these differences, and does not only
correlate with them (with culture, landscape, \ldots)? \textbf{Teach people new
metaphors for time} (e.g.\ English speakers trained to talk about time vertically) and
see if this changes the way they think. Answer: \textbf{yes}.
\end{nbfact}

\subsubsection{Colour}
\slidesrc{Pragmatics: 39--40}

The colour spectrum is continuous, but languages cut it into categories differently.
English has one basic word, \emph{blue}; Russian has two: \textbf{goluboy} (light blue)
and \textbf{siniy} (dark blue). In an experiment (Winawer, Witthoft, Frank, Wade,
Boroditsky, PNAS) Russian speakers were \textbf{faster at distinguishing two blues when
they fell on different sides} of the goluboy/siniy boundary; English speakers showed no
such advantage. The words we have change, measurably, how we perceive.

\subsubsection{Grammatical gender}
\slidesrc{Pragmatics: 41--44}

In German the sun is feminine (\emph{die Sonne}) and the moon masculine (\emph{der
Mond}); in Spanish it is the opposite (\emph{el sol}, \emph{la luna}). Gender is not a
small detail of a few words: it is \textbf{pervasive}. Russian examples with \emph{okno}
(window, neuter), \emph{taburetka} (stool, feminine) and \emph{stul} (chair,
masculine):

\begin{center}
\small
\begin{tabularx}{\linewidth}{@{}>{\raggedright\arraybackslash}p{3.6cm} X@{}}
\toprule
\textbf{where gender shows up} & \textbf{example} \\
\midrule
different phonology & okn\textbf{o} (-o, -e = neuter), taburetk\textbf{a} (-a =
  feminine), stu\textbf{l} (consonant = masculine) \\
different number words & \textbf{odna} taburetka vs.\ \textbf{odin} stul (one stool,
  one chair) \\
different adjective endings & bel\textbf{oye} okno, bel\textbf{aya} taburetka,
  bel\textbf{y} stul (white window/stool/chair) \\
different pronouns and possessives & \textbf{ono} beloye, \textbf{ona} belaya, \textbf{on}
  bely; \textbf{moyo} okno beloye, \textbf{moya} taburetka belaya, \textbf{moy} stul
  bely \\
different verb endings & moyo okno \textbf{bilo} beloye, moya taburetka \textbf{bila}
  belaya, moy stul \textbf{bil} bely (my \ldots\ was white) \\
\bottomrule
\end{tabularx}
\end{center}
To say ``my chair was white'' a Russian speaker must mark the masculine gender of
\emph{chair} on the possessive, the verb and the adjective.

\textbf{Do people really take these genders as meaningful?} Yes: speakers use
grammatical gender when they give voices to animated characters, when they describe
objects, they underestimate hurricanes with female names, and artists personify abstract
entities according to the gender of the word. Michelangelo's statues on the tomb of
Giuliano de' Medici (1526--1531) and of Lorenzo de' Medici represent \emph{Dawn}
(\emph{l'aurora}) and \emph{Night} (\emph{la notte}) as women, \emph{Day} (\emph{il
giorno}) and \emph{Dusk} (\emph{il crepuscolo}) as men: exactly the genders of the
Italian words.

\subsubsection{Agency: who did it?}
\slidesrc{Pragmatics: 45--49}

\emph{``Pat broke the vase.''} English describes even accidents with an agent. Spanish
distinguishes: \emph{Pat rompió el florero} (Pat broke the vase, on purpose) vs.\
\emph{Se rompió el florero} (the vase broke itself, i.e.\ by accident). Language primes
the same contrast in English: \emph{``The match blew out''} vs.\ \emph{``She blew out the
match''}; \emph{``The toast burned''} vs.\ \emph{``She burned the toast''}.

In the experiment, people watched videos of someone breaking things, intentionally or
by accident, and were later asked \textbf{``who did it the first time?''}. For
intentional events everybody remembers the agent equally well; for \textbf{accidental}
events, English speakers remember the agent better than Spanish (or Japanese) speakers,
whose language would normally describe the event without an agent. What a language makes
you say influences what you notice and remember, and, as studies on agentive vs.\
non-agentive descriptions show, also how much blame you assign.

\begin{nbsummary}[language and thought]
\begin{itemize}[nosep]
  \item Newspeak (Orwell): language designed to make thoughts impossible.
  \item Sapir--Whorf: strong (determinism, false) vs.\ weak (influence, supported by
        evidence).
  \item Languages force different information: tense, gender, aspect, evidentiality;
        Navajo verbs of eating. About 7\,160 languages, 44\% endangered.
  \item Time in space: left--right by writing direction; Aymara past in front; Mandarin
        vertical; Kuuk Thaayorre east-to-west; Yupno uphill future. Teaching new
        metaphors changes thinking $\Rightarrow$ causal.
  \item Colour: Russian goluboy/siniy, faster discrimination across the boundary.
  \item Gender: pervasive (Russian), influences descriptions, voices, art (Michelangelo).
  \item Agency: English agentive descriptions; better memory of agents of accidents.
\end{itemize}
\end{nbsummary}
